\subsection{李形式幂级数}

引理 1。设 \(\mathfrak{g}\) 是一个滤李代数（§ 4，第 1 号），\((\mathfrak{g}_{\alpha})_{\alpha \in \mathbf{R}}\) 是其滤子，且 \(\alpha \in \mathbf{R}\)。设 \(\mathbf{P}\) 是 \(n\) 次齐次李多项式，关于不定元 \((\mathbf{T}_i)_{i \in I}\)（§ 2，第 4 号）。则 \(\mathrm{P}((a_i)) \in \mathfrak{g}_{n\alpha}\)，其中 \((a_i)_{i \in I}\) 是由 \(\mathfrak{g}_{\alpha}\) 中元素组成的任意族。

每个次数 \(n \geqslant 2\) 的李多项式都是形如 \([\mathbf{Q}, \mathbf{R}]\) 的项的有限和，其中 \(\mathbf{Q}\) 和 \(\mathbf{R}\) 的次数均为 \(< n\)，且它们的次数之和等于 \(n\)（§ 2，第 7 号，命题 7）。引理由对 \(n\) 作归纳得到。

关于不定元 \((\mathrm{T}_{i})_{i\in\mathbb{I}}\) 的李形式幂级数†（系数在 K 中）是李代数 \(\hat{\mathrm{L}}((\mathrm{T}_{i})_{i\in\mathbb{I}})=\hat{\mathrm{L}}(\mathrm{I})\) 的任一元素。这样的元素 u 可以唯一写成可求和族 \((u_{v})_{v\in\mathbf{N}^{(\mathrm{I})}}\) 的和，其中 \(u_{v}\in\mathrm{L}^{v}(\mathrm{I})\)。

假设 I 是有限集。设 g 是一个完备豪斯多夫滤李代数，满足 \(g = \bigcup_{\alpha > 0} g_{\alpha}\)；令 \(t = (t_i)_{i \in I}\) 为 g 中元素组成的族。

命题 2。满足 \(f_{t} \colon \mathrm{L}(\mathrm{I}) \to \mathfrak{g}\) 的同态 \(f_{t}(\mathrm{T}_{i}) = t_{i}\)（§ 2，第 4 号）可由连续性延拓为唯一的从 \(\hat{f}_{t}\) 到 \(\hat{\mathrm{L}}(\mathrm{I})\) 的连续同态 \(\mathfrak{g}\)。

存在 \(\alpha > 0\)，使得对所有 \(t_i \in \mathfrak{g}_\alpha\) 都有 \(i \in I\)；因此对所有 \(f_t(L^\nu(I)) \subset \mathfrak{g}_{|\nu| \alpha}\)，有 \(\nu\)（引理 1），这蕴含 \(f_t\) 的连续性。

若 \(u \in \widehat{\mathbf{L}}(\mathbf{I})\)，记 \(u((t_i)) = \widehat{f_t}(u)\)。特别地，取 \(\mathfrak{g} = \widehat{\mathbf{L}}(\mathbf{I})\) 时，\(u = u((\mathrm{T}_i))\)；一般情况下，\(u((t_i))\) 称为在李形式幂级数 \(t_i\) 中以 \(T_i\) 代替 \(u((\mathrm{T}_i))\) 所得到的结果。若 \(u = \sum_{v \in \mathbf{N}^{(\mathrm{I})}} u_v\)，其中 \(u_v \in \mathrm{L}^v(\mathrm{I})\)，则族 \((u_v((t_i))_{v \in \mathbf{N}^{(\mathrm{I})}})\) 可求和，并且

\[
u ((t _ {i})) = \sum_ {\nu \in \mathbf {N} ^ {\mathrm{I}}} u _ {\nu} ((t _ {i})).\tag{5}
\]

设 \(\sigma\) 是从 \(\mathfrak{g}\) 到完备豪斯多夫滤李代数 \(\mathfrak{g}'\) 的连续同态，且 \(\mathfrak{g}' = \bigcup_{\alpha > 0} \mathfrak{g}_{\alpha}'\)。对于 \(\boldsymbol{t} = (t_i)_{i \in I}\) 中元素组成的每个有限族 \(\mathfrak{g}\) 以及所有 \(u \in \hat{\mathbf{L}}(\mathbf{I})\)，

\[
\sigma (u ((t _ {i}))) = u ((\sigma (t _ {i}))),\tag{6}
\]

因为同态 \(\sigma \circ f_t\) 是连续的，并且将 \(T_i\) 映到 \(\sigma(t_i)\)（\(i \in I\)）。

令 \(\boldsymbol{u} = (u_j)_{j \in \mathcal{J}}\) 为 \(\hat{\mathrm{L}}(\mathrm{I})\) 中元素组成的有限族，并令 \(v \in \hat{\mathrm{L}}(\mathrm{J})\)；通过用 \(u_j\) 代替 \(T_j\) 并代入 \(v\)，得到元素 \(w = v((u_j)_{j \in \mathcal{J}})\)，它属于 \(\hat{\mathrm{L}}(\mathrm{I})\)，记为 \(v \circ \boldsymbol{u}\)。于是

\[
w ((t _ {i}) _ {i \in I}) = v ((u _ {j} ((t _ {i}) _ {i \in I})) _ {j \in J})\tag{7}
\]

对于 g 中元素组成的每个有限族 \(\boldsymbol{t} = (t_{i})_{i \in I}\) 都成立，这是对等式 \(\hat{f}_{t}\) 施加连续同态 \(w = v((u_{j})_{j \in J})\) 即可看出。

令 \(u = \sum_{\nu \in \mathbb{N}^{\mathbf{I}}} u_{\nu} \in \hat{\mathrm{L}}(\mathrm{I})\)，其中 \(u_{\nu} \in \mathrm{L}^{\nu}(\mathrm{I})\)。从 \(\tilde{u}\colon (t_i)\mapsto u((t_i))\) 到 \(\mathfrak{g}^{\mathbf{I}}\) 的映射 \(\mathfrak{g}\) 是连续的：因为在每个 \(\mathfrak{g}_{\alpha}\) 的开集 \(\alpha > 0\) 中，族 \(\tilde{u}_{\nu}\) 都一致可求和，只需证明每个 \(\tilde{u}_{\nu}\) 连续即可，而这由对 \(|\nu|\) 作归纳立即得到。

